\documentclass[10pt,a4paper,openany,twoside]{book}

\usepackage{layout_cours}


\begin{document}

    \titre{Calcul algébrique} % titre du chapitre
	\classe{Seconde} % classe

    \pagenumbering{Alph}
    \enteteperso
    \section*{Bordas 2023}
    \includegraphics[width=.6\textwidth]{Bordas/Bordas_Indice_2nde_2023_p_105_Z6.jpg}

    \includegraphics[width=.6\textwidth]{Bordas/Bordas_Indice_2nde_2023_p_106_Z2.jpg}

    \includegraphics[width=.6\textwidth]{Bordas/Bordas_Indice_2nde_2023_p_106_Z7.jpg}

\newpage

    \thispagestyle{premierepage}
    \pagenumbering{arabic} %need to reset the pagenumbering to call it as a new sheme
    \pagenumbering{Alph}
    \enteteperso
    \section*{Bordas 2023}

    \includegraphics[width=.6\textwidth]{Bordas/Bordas_Indice_2nde_2023_p_105_Z6.jpg}

    \includegraphics[width=.6\textwidth]{Bordas/Bordas_Indice_2nde_2023_p_106_Z2.jpg}

    \includegraphics[width=.6\textwidth]{Bordas/Bordas_Indice_2nde_2023_p_106_Z7.jpg}

\newpage
\pagenumbering{arabic}
\enteteperso
\section*{Bordas 2023}
\sectionexo{Développer avec la simple et la double distributivité}
\begin{bookexos}{125-128 p.105}\end{bookexos}
\sectionexo{Développer avec les identités remarquables}
\begin{bookexos}{131-133 p.106}\end{bookexos}
\sectionexo{Factoriser}
\begin{bookexos}{138-141 p.106}\end{bookexos}
\sectionexo{Factoriser avec les identités remarquables}
\begin{bookexos}{142-143 p.106}\end{bookexos}

\newpage
\section*{Corrigé}
\begin{bookexo}{125 p.105}
$\begin{aligned}
    x-4(3x+2) &= -11x+8\\
    (x-2)(x-3) &=x^2-5x+6\\
    6-3(2x-1) &=-6x+9\\
    (5x-4)(3x+2) &= 15x^2-2x-8
\end{aligned}$
\end{bookexo}

\begin{bookexo}{126 p.105}
$\begin{aligned}
    5x-(x+1)(6x-2)
        &=5x-[6x^2+6x-2x-2]\\
        &=5x-6x^2-4x+2\\
        &=-6x^2+x+2
\end{aligned}$

$\begin{aligned}
    5(\frac{2}{5}x-3)(x-1)
        &= (2x-15)(x-1)\\
        &= 2x^2-2x-15x+15 \\
        &= 2x^2-17x+15
\end{aligned}$
\end{bookexo}

\begin{bookexo}{127 p.105}
$\begin{aligned}
    2-(x+3)(2x+1)
        &= 2-[2x^2+x+6x+1]\\
        &= -2x^2-7x-1
\end{aligned}$

$\begin{aligned}
    x(x-1)(x+2)
        &=  x[x^2+2x-x-2]\\
        &=  x[x^2+x-2]\\
        &=  x^3+x^2-2x
\end{aligned}$
\end{bookexo}

\begin{bookexo}{128 p.105}

$\begin{aligned}
    3(x-2)(2x+5)
        &=  3(2x^2+5x-4x-10)\\
        &=  6x^2+3x-30
\end{aligned}$

$\begin{aligned}
    (x-1)(x^2+x+1)
        &=  x^3+&x^2+&x\\
        &       &-x^2-&x-1\\
        &=  x^3-1
\end{aligned}$
\end{bookexo}

\begin{bookexo}{131 p.106}
$\begin{aligned}
    (x+3)^2&=x^2+6x+9 \\
    (x-4)(x+4)&=x^2-16 \\
    (x-5)^2&=x^2-10x+25
\end{aligned}$
\end{bookexo}

\begin{bookexo}{132 p.106}
$\begin{aligned}
    (3x+5)^2 &= 9x^2+30x+25 \\
    (2x-1)(2x+1) &= 4x^2-1 \\
    (1-2x)^2 &= -4x^2-4x+1
\end{aligned}$
\end{bookexo}

\begin{bookexo}{133 p.106}
$\begin{aligned}
    (x-4)^2+(x+2)^2
        &=&x^2&-8x&+16\\
        & &+x^2&+4x&+4\\
        &=&2x^2&-4x&+20
\end{aligned}$

$\begin{aligned}
    (x+5)^2-(2x-4)^2
        &=& &x^2&+10x&+25\\
        &&  &-4x^2&+16&x-16\\
        &=& &3x^2&+26x&+9
\end{aligned}$
\end{bookexo}

\begin{bookexo}{138 p.106}
$\begin{aligned}
4(x-3)+x(x-3)&=(4+x)(x-3)=(x+4)(x-3) \\
x^2-4x &= x(x-4)
\end{aligned}$
\end{bookexo}

\begin{bookexo}{139 p.106}
$\begin{aligned}
    (x+2)(x-1)+(3x+5)(x+2)
        &= (x+2)(x-1+3x+5)\\
        &= (x+2)(4x+4)\\
        &= 4(x+2)(x+1)
\end{aligned}$

$\begin{aligned}
    7(5-x)-(2x+3)(5-x)
        &= (7-(2x+3))(5-x) \\
        &= (-2x+4)(5-x) \\
        &= 2(-x+2)(-x+5)
\end{aligned}$
\end{bookexo}

\begin{bookexo}{140 p.106}
$\begin{aligned}
    (3x+1)(x+4)-(2x-5)(x+4)
        &= (3x+1-(2x-5))(x+4) \\
        &= (x+6)(x+4)
\end{aligned}$

$\begin{aligned}
    (7x-3)^2-x(7x-3)
        &= (7x-3)(7x-3-x) \\
        &= (7x-3)(6x-3)\\
        &= 3(7x-3)(2x-1)
\end{aligned}$
\end{bookexo}



% $\begin{aligned}

%         &=  \\
%         &=  \\
%         &=  \\
%         &=
% \end{aligned}$

\begin{bookexo}{141 p.106}
    $\begin{aligned}
        A(x)
        &=(2-x)[1-(2-x)]\\
        &=(2-x)(-1+x)\\
        &=(-x+2)(x-1)
    \end{aligned}$

    $\begin{aligned}
    B(x)
        &=(x+1)^2-5(x+1) \\
        &=(x+1)(x+1-5) \\
        &=(x+1)(x-4)
    \end{aligned}$
\end{bookexo}

\begin{bookexo}{142 p.106}
$\begin{aligned}
    A(x)
        &=(x-2)^2-36 \\
        &=(x-2-6)(x-2+6) \\
        &=(x-8)(x+4)
\end{aligned}$

$\begin{aligned}
    B(x)
        &= 1-(2x+3)^2 \\
        &= [1-(2x+3)][1+(2x+3)] \\
        &= (-2x-2)(2x+4) \\
        &= -4(x+1)(x+2)
\end{aligned}$
\end{bookexo}

\begin{bookexo}{143 p.106}
$\begin{aligned}
    A(x)
        &= 16-(x-5)^2 \\
        &= [4-(x-5)][4+(x-5)] \\
        &= (-x+9)(x-1)
\end{aligned}$

$\begin{aligned}
    B(x)
        &=(3x-2)^2-(x+3)^2 \\
        &=[(3x-2)-(x+3)][(3x-2)+(x+3)] \\
        &= (2x-5)(4x+1)
\end{aligned}$
\end{bookexo}



\end{document}